% sec:rk-adaptive — Adaptive Step-Size Control
\section{Adaptive Step-Size Control}
\label{sec:rk-adaptive}

A photon cruising through nearly flat spacetime far from the black hole
needs nothing like the step resolution required during a close pass
around the photon sphere.  A fixed step size forces an impossible
choice: either waste orders of magnitude in compute on the gentle
segments, or risk corrupting the trajectory on the hard ones.  The
error estimate $\delta y_{n+1}$ from \cref{sec:rk-dormand-prince}
lets the integrator resolve this tension by adjusting its step size on
the fly --- but three decisions determine how: how to measure the size
of the error, when to accept or reject a step, and how to choose the
next step size.
\xrefnote{This error norm was previewed in \cref{sec:geo-raytracing};
here we derive the step-size controller it feeds into.}

\paragraph{The scaled error norm.}

The raw error vector $\delta y_{n+1}$ has the same dimension as the
state $y$ --- for the geodesic integrator, eight components comprising
four position coordinates and four momentum components.  Comparing
these components requires scaling: a position error of $10^{-8}$ in
geometric units and a momentum error of $10^{-8}$ may have very
different physical significance depending on the magnitudes of the
state variables.  The standard approach scales each component by a
tolerance that blends absolute and relative thresholds:
%
\begin{equation}
  \mathrm{err} = \max_{i} \frac{\abs{\delta y_i}}
    {\epsilon_a + \epsilon_r\,\abs{y_{n+1,i}}} \,.
  \label{eq:rk-error-norm}
\end{equation}
%
The absolute tolerance $\epsilon_a$ prevents division-by-zero issues
when $y_i$ is near zero, and the relative tolerance $\epsilon_r$ tracks
the scale of each component as it evolves.  The denominator uses the
proposed state $y_{n+1}$ rather than the current state $y_n$, so that
the scaling reflects the regime the integrator is entering.  The
maximum over all components ensures that no single component violates
the tolerance: when $\mathrm{err} = 1$, the worst-offending component
sits exactly at the tolerance boundary; values below one mean every
component is within tolerance, and values above one signal that at
least one component has exceeded it.
\physnote{For the geodesic integrator, $y = (x^\mu, p_\mu)$.  The
tolerances $\epsilon_a$ and $\epsilon_r$ apply uniformly to all eight
components; the relative tolerance handles scale differences between
large-magnitude components, while $\epsilon_a$ protects components
passing through zero.}

\paragraph{Accept, reject, resize.}

The scaled error norm is constructed so that $\mathrm{err} \leq 1$
means every component of the error is within tolerance.  The
controller's logic is then straightforward:
%
\begin{equation}
  \text{step accepted} \;\;\Longleftrightarrow\;\; \mathrm{err} \leq 1 \,.
  \label{eq:rk-accept}
\end{equation}
%
If accepted, the integrator advances the state to $y_{n+1}$ and
(via FSAL) passes $k_7$ as the first stage of the next step.  If
rejected, the state remains at $y_n$, the existing $k_1$ is retained,
and the integrator retries with a smaller step size.

The key question is how to choose the next step size $h_{\text{new}}$.
The error estimate scales as $\mathrm{err} \propto h^5$ (since
$\delta y_{n+1} = \order{h^5}$ from \cref{eq:embedded-error}, and the
tolerance scale in the denominator of \cref{eq:rk-error-norm} is
independent of $h$, so the max-norm inherits the fifth-order scaling).
Assuming the local error coefficient is roughly constant between
successive steps gives
$\mathrm{err}_{\text{new}} / \mathrm{err}
  \approx (h_{\text{new}} / h)^5$.
Setting $\mathrm{err}_{\text{new}} = 1$ and solving for
$h_{\text{new}}$ yields
$h_{\text{new}} = h \cdot \mathrm{err}^{-1/5}$.
Multiplying by a safety factor $S < 1$ gives the working formula:
%
\begin{equation}
  h_{\text{new}} = h \cdot S \cdot \mathrm{err}^{-1/5} \,,
  \label{eq:rk-step-control}
\end{equation}
%
where the exponent $-1/5$ follows directly from the fifth-order scaling
of the error estimate.  The safety factor (typically $S = 0.9$)
prevents the controller from stepping right up to the acceptance
threshold, which would cause frequent rejections due to higher-order
terms neglected in the asymptotic error estimate.
\coderef{Geodesic}{Solver}

\paragraph{Growth clamps.}

The raw scaling factor $S \cdot \mathrm{err}^{-1/5}$ can be
arbitrarily large when $\mathrm{err}$ is very small (the step was much
more accurate than needed) or dangerously aggressive when
$\mathrm{err}$ is only slightly above $1$.  The controller clamps the
factor to a bounded range:
%
\begin{equation}
  h_{\text{new}} = h \cdot
  \operatorname{clamp}\!\bigl(S \cdot \mathrm{err}^{-1/5},\;
    s_{\min},\; s_{\max}\bigr) \,,
  \label{eq:rk-clamp}
\end{equation}
%
where $\operatorname{clamp}(x, a, b) = \min(b, \max(a, x))$.  The
codebase uses $s_{\min} = 0.2$ and $s_{\max} = 5.0$: the step size is
never shrunk by more than a factor of five or grown by more than a
factor of five in a single step.  Without the upper clamp, a photon in
nearly flat spacetime could jump to an enormous step that overshoots
its approach to the black hole; without the lower clamp, a single
noisy step near the photon sphere could collapse $h$ to a value from
which the integrator never recovers.
\implnote{The \texttt{SolverConfig} record in
\codemodule{Geodesic.Solver} exposes \texttt{scSafety},
\texttt{scMinScale}, and \texttt{scMaxScale} as configurable
parameters, with defaults $S = 0.9$, $s_{\min} = 0.2$,
$s_{\max} = 5.0$.}

\paragraph{Rejected-step shrinkage.}

When a step is rejected ($\mathrm{err} > 1$), the controller must
shrink $h$ and retry.  Using the same exponent $-1/5$ would produce a
moderate shrinkage, but the error estimate may be less reliable when
it is large: the $\order{h^5}$ asymptotic regime requires $h$ small
enough that higher-order terms are negligible, and a step that failed
the tolerance test may not be in that regime.  A common heuristic is
to use a more conservative exponent for rejected steps:
%
\begin{equation}
  h_{\text{new}} = h \cdot
  \max\bigl(s_{\min},\;
    S \cdot \mathrm{err}^{-1/4}\bigr) \,,
  \qquad (\text{rejected step}) \,,
  \label{eq:rk-reject}
\end{equation}
%
with exponent $-1/4$ rather than $-1/5$.  Using a larger-magnitude
exponent is equivalent to assuming the error grows faster than $h^5$
in this regime, producing a more conservative step size as a margin of
safety and reducing the chance of a second consecutive rejection.  The
$s_{\min}$ floor still applies, ensuring the step size cannot collapse
below $0.2\,h$.  The upper clamp $s_{\max}$ is omitted because
$\mathrm{err} > 1$ and $S < 1$ guarantee the scaling factor is less
than one.

\warnnote{A rejected step does not advance the state or the step
counter.  The integrator retries from the same $(t_n, y_n)$ with the
shrunken step size and the same $k_1$ (valid by the FSAL argument of
\cref{sec:rk-fsal}).}

\paragraph{Default parameters.}

The codebase's default solver configuration uses tight tolerances
appropriate for a reference implementation:
%
\begin{equation}
  \begin{gathered}
  \epsilon_a = 10^{-10} \,, \quad
  \epsilon_r = 10^{-10} \,, \quad
  S = 0.9 \,, \quad
  h_0 = 0.5 \,, \\[4pt]
  s_{\min} = 0.2 \,, \quad
  s_{\max} = 5.0 \,.
  \end{gathered}
  \label{eq:rk-defaults}
\end{equation}
%
The initial step size $h_0 = 0.5$ (in affine-parameter units) is
deliberately generous: the controller will shrink it rapidly during the
first few steps as it discovers the appropriate scale for the
trajectory.  The tight tolerances $\epsilon_a = \epsilon_r = 10^{-10}$
ensure that the Hamiltonian constraint $\mathcal{H} = 0$ is preserved
to near machine precision over long integrations, as verified by the
conservation diagnostics of \cref{sec:geo-conservation}.

In practice, the adaptive controller reduces the cost of a ray trace
by one to two orders of magnitude compared with a fixed-step integrator
at the same accuracy.  A typical photon might cruise at $h \sim 10$
through the asymptotically flat region but shrink to $h \sim 10^{-3}$
during a close encounter with the photon sphere --- a dynamic range
that no single fixed step size can efficiently span.
